%The conclusions of this work and a plan of action for the following months of work are presented in this chapter. The steps are described and distributed
%through a project schedule.

%\section{Future Works}


%\section{Conclusions}
% In this work, a revision of the concepts related to the design of a control method based on images for a two-degrees-of-freedom gimbal was presented. With the concepts presented, it is possible to understand how the proposed solution will be designed. Having the right tools surely facilitates the development process, as proven by the results of YOLOv8, which gives enough information to understand the opportunities for improvement. With the Kalman filter algorithm, it is possible to estimate the behavior of a system using indirectly measured data, which is useful for object-tracking applications. With the study presented in Section \ref{secaogimbal}, we see that modeling the dynamics of a gimbal is not trivial, but some considerations can be made to simplify this process. Also, with the help of a project schedule tool such as the Gantt chart, the visualization of the next steps becomes easier, ensuring the completion of all tasks at the appropriate time.

This chapter consolidates the main conclusions drawn from the development and validation of the image-based gimbal control system. It also outlines potential avenues for future research, aiming to further enhance the system's performance and applicability.

\section{Conclusions}
This Master's thesis successfully developed and validated a comprehensive control solution for real-time object tracking in images, utilizing a two-degrees-of-freedom (2-DOF) gimbal system. The work addressed the inherent challenges associated with complex gimbal dynamics, including nonlinearities and communication delays, by proposing an integrated architecture.

The core of the proposed solution involved the synergistic integration of a state-of-the-art object detection algorithm (YOLOv8), a state estimator (Kalman Filter) for signal smoothing, and a Proportional-Integral-Derivative (PID) control system. The control actuation was applied to a detailed Simulink model of a 2-DOF gimbal, responsible for adjusting the camera's orientation to maintain the object centered within the field of view.

Several key findings and achievements emerged from this study. A robust end-to-end system architecture was established, seamlessly connecting visual perception (YOLOv8), data communication (UDP protocol), state estimation (Kalman Filter), and motion control (PID controller) with the electromechanical dynamics of the gimbal. This integration was validated in a unified Simulink environment, demonstrating a fundamental step towards autonomous tracking systems. Regarding object detection performance, the YOLOv8 algorithm demonstrated its capability in real-time object detection, providing the necessary pixel coordinates for tracking. While general performance metrics (mAP@50 and mAP@50-95) showed promising trends during training, a class imbalance within the VisDrone2019 dataset was identified as a factor influencing detection accuracy across different object categories.

System latency management was a critical aspect, and the overall system latency, encompassing camera delay, YOLO processing time, UDP communication, and the Controller and Actuator delays, resulted in approximately 124 milliseconds. The PID controller, designed using MATLAB's \texttt{pidtune()} function and refined through fine-tuning, demonstrated its ability to stabilize the gimbal's rotational axes. For the elevation axis, the step response showed stable behavior with an initial overshoot attributed to the transient effects of communication startup and noise. The system effectively tracked the zero reference (center of the image) with an approximate settling time of 5 seconds, exhibiting an underdamped response. The azimuth axis exhibited similar dynamic characteristics to the elevation axis, as expected due to the simplified model's assumptions of analogous dynamics. It achieved a lower settling time but presented a more oscillatory behavior, indicating areas for potential improvement. Finally, the gimbal model incorporated critical nonlinearities such as viscous and Coulomb frictions and considered the rotational inertia of the motor and encoder, enhancing the realism of the simulation environment. The inclusion of communication delay via Padé approximation further contributed to a more faithful evaluation under conditions emulating the physical world.

In conclusion, this work successfully demonstrated the viability of an image-based control system for a 2-DOF gimbal, providing a stable and responsive platform for object tracking. The integrated approach, coupled with careful modeling of system dynamics and real-time considerations, lays a solid foundation for future advancements in autonomous visual tracking systems.

\section{Future Works}
Building upon the achievements and insights gained from this master's thesis, several avenues for future research are identified to further enhance the performance, robustness, and real-world applicability of the image-based gimbal control system.

One area for future exploration involves advanced control strategies. This could include investigating the implementation of advanced nonlinear control techniques, such as Sliding Mode Control (SMC) \cite{Rahmat2016Sliding}, since this method is inherently better suited to handle the complex nonlinearities (e.g., Coulomb friction, actuator saturation, torque ripple) present in gimbal dynamics, potentially leading to improved transient response (reduced overshoot and oscillations) and enhanced disturbance rejection. Furthermore, adaptive control algorithms, such as Fuzzy logic controller \cite{Galal2009Fuzzy}, could be investigated to allow the controller to adjust its parameters in real-time in response to varying operating conditions or unmodeled dynamics, thereby increasing system robustness. Developing explicit compensation strategies for known disturbances, such as dynamic mass unbalance and torque ripple, could further improve tracking precision and reduce steady-state errors. Finally, implementing control strategies specifically designed to actively compensate for the measured communication and processing delays, rather than merely tolerating them, could involve predictive control elements or Smith predictors.

Another significant area for future work lies in object detection and state estimation enhancement. This includes addressing the identified dataset imbalance in YOLO training through resampling techniques or data augmentation strategies, and further optimizing YOLOv8 hyperparameters to improve detection accuracy across all classes, especially those with fewer instances. Exploring and comparing the performance of other state-of-the-art object detection algorithms (e.g., Faster R-CNN, SSD variants) for their suitability in real-time gimbal control applications, considering their trade-offs between accuracy and computational cost, is also a valuable direction. Additionally, extending the Kalman filter-based state estimator to track multiple objects simultaneously would require more sophisticated data association and state management techniques. Enhancing capabilities to estimate a target's distance and velocity, particularly when associated with ballistics, expands its applicability to military systems.

Gimbal model refinement presents further opportunities. This could involve incorporating more granular physical phenomena into the Simulink model, such as detailed bearing dynamics, structural flexibility of the gimbal arms, and more precise modeling of cross-coupling effects that were simplified in the current iteration. Integrating and compensating for the dynamic angular motion of the base (e.g., aircraft maneuvers), which was simplified by the fixed-base assumption in the current model, is critical for real-world applications.

Experimental validation is a crucial next step. This involves transitioning the validated control solution from the simulation environment to a physical 2-DOF gimbal prototype. This step is essential for real-world testing, allowing for the identification and mitigation of unmodeled dynamics and practical challenges not captured in simulation. Implementing Real-time Hardware-in-the-Loop (HIL) simulations to test the controller with actual hardware components (e.g., motors, encoders) connected to the simulated gimbal model would provide a more realistic validation environment before full physical deployment.

Finally, computational optimization is an important consideration. Optimizing the computational efficiency of the object detection and control algorithms for deployment on embedded systems or edge computing devices is typical in UAV applications.

These future works collectively aim to enhance the system's robustness, precision, and practical applicability, paving the way for more sophisticated and reliable image-based gimbal control systems in diverse engineering domains.